% C057 | 双曲线焦点分弦定理 | 2026-10-26 修订
% 本文件为内容片段，依赖主文档已定义的 statementbox、CStatement 等样式。
\begin{statementbox}
  \textbf{\textcolor{CStatement}{条件（先判交点，再套结论）}}
  \begin{description}[style=nextline,leftmargin=2.8em,labelsep=0.5em,itemsep=0.35em]
    \item[\textbf{条件1（双曲线）：}] $\displaystyle\frac{x^2}{a^2}-\frac{y^2}{b^2}=1$，其中 $a>0,b>0$，$c^2=a^2+b^2$，焦点 $F(\pm c,0)$。
    \item[\textbf{条件2（过焦点的直线）：}] 直线 $\ell$ 经过选定焦点 $F$，且与双曲线有\textbf{两个不同的有限交点} $A,B$。此处“焦点弦”指\textbf{弦所在直线过焦点}，不要求 $F$ 在线段 $AB$ 内。
    \item[\textbf{条件3（分支）：}] 用 $x_A,x_B$ 的正负判断同支或异支；过右焦点的同支两交点必在右支，过左焦点的同支两交点必在左支。
  \end{description}

  \medskip
  \textbf{\textcolor{CStatement}{核心结论}}
  \begin{description}[style=nextline,leftmargin=2.8em,labelsep=0.5em,itemsep=0.45em]
    \item[\textbf{结论一（同支：倒数相加）：}]
      当 $A,B$ 同支时，$F$ 位于线段 $AB$ 内，且
      \[
        \boxed{\frac1{|AF|}+\frac1{|BF|}=\frac{2a}{b^2}}.
      \]
    \item[\textbf{结论二（异支：倒数相减）：}]
      当 $A,B$ 异支时，$F$ 位于线段 $AB$ 的延长线上，且
      \[
        \boxed{\left|\frac1{|AF|}-\frac1{|BF|}\right|=\frac{2a}{b^2}}.
      \]
      \textbf{更强的有序式：} 若 $F=F_2(c,0)$，$A$ 在右支、$B$ 在左支，必有 $|AF|<|BF|$，故
      \[
        \boxed{\frac1{|AF|}-\frac1{|BF|}=\frac{2a}{b^2}}.
      \]
      左焦点时将“右支、左支”互换。
    \item[\textbf{统一常数（不用背两份）：}]
      设半通径参数 $p=\dfrac{b^2}{a}$、离心率 $e=\dfrac ca>1$，则
      \[
        \boxed{\frac{2a}{b^2}=\frac2p=\frac{2}{a(e^2-1)}}.
      \]
  \end{description}

  \medskip
  \textbf{\textcolor{CStatement}{范围提醒：}}
  斜率 $k=\pm\dfrac ba$ 的过焦点直线与一条渐近线平行，仅有一个有限交点，\textbf{不构成}上述两交点情形；竖直线 $x=\pm c$ 则有两个同支交点，\textbf{可以使用}结论。
\end{statementbox}
